% Slug: ch34-tiers
% Source: chapters/ch34.tex:30 (tikzpicture env start line)
% Original label: \label{fig:ch34-tiers}
% Original caption (LaTeX source):
%   Tiered execution moves hot code upward as counters grow and returns through deoptimization when assumptions fail, so profiling pays only where it matters.
% Original caption (plain text):
%   Tiered execution moves hot code upward as counters grow and returns through deoptimization when assumptions fail, so profiling pays only where it matters.
% Required TikZ setup: \usepackage{tikz} + \usetikzlibrary{arrows.meta, positioning, shapes.geometric, fit} (loaded by support/compiler-pseudocode.sty).
% Shared styles used in this picture: `basic block', `cfg edge'.
%   Defined in support/compiler-pseudocode.sty: `basic block', `cfg edge', `back edge', `block label'.
% Usage: % Slug: ch34-tiers
% Source: chapters/ch34.tex:30 (tikzpicture env start line)
% Original label: \label{fig:ch34-tiers}
% Original caption (LaTeX source):
%   Tiered execution moves hot code upward as counters grow and returns through deoptimization when assumptions fail, so profiling pays only where it matters.
% Original caption (plain text):
%   Tiered execution moves hot code upward as counters grow and returns through deoptimization when assumptions fail, so profiling pays only where it matters.
% Required TikZ setup: \usepackage{tikz} + \usetikzlibrary{arrows.meta, positioning, shapes.geometric, fit} (loaded by support/compiler-pseudocode.sty).
% Shared styles used in this picture: `basic block', `cfg edge'.
%   Defined in support/compiler-pseudocode.sty: `basic block', `cfg edge', `back edge', `block label'.
% Usage: % Slug: ch34-tiers
% Source: chapters/ch34.tex:30 (tikzpicture env start line)
% Original label: \label{fig:ch34-tiers}
% Original caption (LaTeX source):
%   Tiered execution moves hot code upward as counters grow and returns through deoptimization when assumptions fail, so profiling pays only where it matters.
% Original caption (plain text):
%   Tiered execution moves hot code upward as counters grow and returns through deoptimization when assumptions fail, so profiling pays only where it matters.
% Required TikZ setup: \usepackage{tikz} + \usetikzlibrary{arrows.meta, positioning, shapes.geometric, fit} (loaded by support/compiler-pseudocode.sty).
% Shared styles used in this picture: `basic block', `cfg edge'.
%   Defined in support/compiler-pseudocode.sty: `basic block', `cfg edge', `back edge', `block label'.
% Usage: % Slug: ch34-tiers
% Source: chapters/ch34.tex:30 (tikzpicture env start line)
% Original label: \label{fig:ch34-tiers}
% Original caption (LaTeX source):
%   Tiered execution moves hot code upward as counters grow and returns through deoptimization when assumptions fail, so profiling pays only where it matters.
% Original caption (plain text):
%   Tiered execution moves hot code upward as counters grow and returns through deoptimization when assumptions fail, so profiling pays only where it matters.
% Required TikZ setup: \usepackage{tikz} + \usetikzlibrary{arrows.meta, positioning, shapes.geometric, fit} (loaded by support/compiler-pseudocode.sty).
% Shared styles used in this picture: `basic block', `cfg edge'.
%   Defined in support/compiler-pseudocode.sty: `basic block', `cfg edge', `back edge', `block label'.
% Usage: \input{graphs-code/ch34-tiers.tikz.tex} inside a figure environment
%   (wrap with \begin{figure} ... \caption{...} \label{fig:ch34-tiers} ... \end{figure}).
% Note: tikzpicture body below is verbatim from the source; do not reformat.
\begin{tikzpicture}[node distance=1.0cm and 1.4cm]
\node[basic block] (interp) {interpreter\\ profile};
\node[basic block, right=of interp] (base) {baseline JIT\\ fast, modest code};
\node[basic block, right=of base] (opt) {optimizing JIT\\ slow, best code};
\node[basic block, below=of base] (deopt) {deoptimize\\ back to lower tier};
\draw[cfg edge] (interp) -- (base);
\draw[cfg edge] (base) -- (opt);
\draw[cfg edge] (opt) -- (deopt);
\draw[cfg edge] (deopt) -- (interp);
\end{tikzpicture}
 inside a figure environment
%   (wrap with \begin{figure} ... \caption{...} \label{fig:ch34-tiers} ... \end{figure}).
% Note: tikzpicture body below is verbatim from the source; do not reformat.
\begin{tikzpicture}[node distance=1.0cm and 1.4cm]
\node[basic block] (interp) {interpreter\\ profile};
\node[basic block, right=of interp] (base) {baseline JIT\\ fast, modest code};
\node[basic block, right=of base] (opt) {optimizing JIT\\ slow, best code};
\node[basic block, below=of base] (deopt) {deoptimize\\ back to lower tier};
\draw[cfg edge] (interp) -- (base);
\draw[cfg edge] (base) -- (opt);
\draw[cfg edge] (opt) -- (deopt);
\draw[cfg edge] (deopt) -- (interp);
\end{tikzpicture}
 inside a figure environment
%   (wrap with \begin{figure} ... \caption{...} \label{fig:ch34-tiers} ... \end{figure}).
% Note: tikzpicture body below is verbatim from the source; do not reformat.
\begin{tikzpicture}[node distance=1.0cm and 1.4cm]
\node[basic block] (interp) {interpreter\\ profile};
\node[basic block, right=of interp] (base) {baseline JIT\\ fast, modest code};
\node[basic block, right=of base] (opt) {optimizing JIT\\ slow, best code};
\node[basic block, below=of base] (deopt) {deoptimize\\ back to lower tier};
\draw[cfg edge] (interp) -- (base);
\draw[cfg edge] (base) -- (opt);
\draw[cfg edge] (opt) -- (deopt);
\draw[cfg edge] (deopt) -- (interp);
\end{tikzpicture}
 inside a figure environment
%   (wrap with \begin{figure} ... \caption{...} \label{fig:ch34-tiers} ... \end{figure}).
% Note: tikzpicture body below is verbatim from the source; do not reformat.
\begin{tikzpicture}[node distance=1.0cm and 1.4cm]
\node[basic block] (interp) {interpreter\\ profile};
\node[basic block, right=of interp] (base) {baseline JIT\\ fast, modest code};
\node[basic block, right=of base] (opt) {optimizing JIT\\ slow, best code};
\node[basic block, below=of base] (deopt) {deoptimize\\ back to lower tier};
\draw[cfg edge] (interp) -- (base);
\draw[cfg edge] (base) -- (opt);
\draw[cfg edge] (opt) -- (deopt);
\draw[cfg edge] (deopt) -- (interp);
\end{tikzpicture}
